\documentclass[11pt,letterpaper]{article}

\usepackage[spanish,es-noshorthands]{babel}
\usepackage{fontspec}
\setmainfont{Source Serif Pro}
\usepackage{microtype}
\usepackage{geometry}
\usepackage{booktabs}
\usepackage{tabularx}
\usepackage{enumitem}
\usepackage{xcolor}
\usepackage{graphicx}
\usepackage{csquotes}
\usepackage{amsmath}
\usepackage{siunitx}
\usepackage{tikz}
\usetikzlibrary{arrows.meta,positioning,shapes.geometric,fit,backgrounds}
\usepackage{xurl}
\usepackage[hidelinks]{hyperref}
\usepackage[backend=biber,style=apa,sorting=nyt]{biblatex}
\DeclareLanguageMapping{spanish}{spanish-apa}

\geometry{margin=2.3cm}
\setlength{\emergencystretch}{1em}
\setlength{\parindent}{0pt}
\setlength{\parskip}{0.55em}
\setlist{nosep}
\addbibresource{referencias.bib}
\AtBeginBibliography{\sloppy}

\sisetup{
  output-decimal-marker = {,},
  group-separator = {.},
  group-minimum-digits = 4,
  detect-all
}

\title{\vspace{-1.2cm}\textbf{Patentes y globalización: colaborar no es lo mismo
que controlar}\\
\large Por qué considero que la patente es un mal necesario}
\author{Carlos Luis Rojas Aragonés\\
\small Gestión del Desarrollo Tecnológico}
\date{\small 7 de agosto de 2026}

\begin{document}

\maketitle
\vspace{-0.8em}

\section*{Qué mide realmente la tendencia}

Que la patente sirva como indicador de globalización no es una intuición: es una
propiedad del documento, que registra dónde se inventó, quién es el titular y en
qué mercados se busca protección \parencite{guellec2001,oecd2009}. Eso permite
separar tres fenómenos que el enunciado del foro agrupa en uno solo, siguiendo
la taxonomía de \textcite{archibugi1995}: explotación global de tecnología
nacional, colaboración tecnológica global y generación global de innovación.

La distinción importa. La \textbf{co-invención internacional} es cooperación:
equipos de países distintos firman juntos la misma invención. El \textbf{control
transfronterizo} es propiedad: quién termina siendo titular, con frecuencia una
matriz extranjera sobre una invención generada en un centro de I+D local. Lo
primero reparte capacidades; lo segundo reparte rentas. Pueden ir juntos, pero
no son el mismo hecho ni merecen el mismo juicio. La escala del sistema ayuda a
dimensionar lo que está en juego: en 2024 se presentaron \num{3.7} millones de
solicitudes de patente en el mundo, un \SI{4.9}{\percent} más que en 2023, y la
oficina china recibió \num{1.8} millones, el \SI{49.1}{\percent} del total
mundial frente al \SI{34.6}{\percent} de 2014 \parencite{wipo2025}.

\begin{figure}[ht]
  \centering
  \begin{tikzpicture}[xscale=1.02,yscale=0.60]
    \draw[-{Stealth[length=2mm]}] (0,0) -- (10.9,0);
    \node[font=\scriptsize] at (10.55,-0.85) {Año};
    \draw[-{Stealth[length=2mm]}] (0,0) -- (0,7.1)
      node[above right=-1mm and -1.4cm, font=\scriptsize]
      {Solicitudes PCT (miles)};
    \foreach \y/\lab in {0/0, 1/50, 2/100, 3/150, 4/200, 5/250, 6/300} {
      \draw[gray!25] (0,\y) -- (10.4,\y);
      \node[left, font=\scriptsize] at (-0.08,\y) {\lab};
    }
    \foreach \x/\lab in {0/2000, 2/2005, 4/2010, 6/2015, 8/2020, 10/2025} {
      \draw (\x,0) -- (\x,-0.18);
      \node[below, font=\scriptsize] at (\x,-0.18) {\lab};
    }
    \fill[orange!13] (8,0) rectangle (9.6,6.3);
    \node[font=\scriptsize, color=orange!60!black] at (8.8,6.75)
      {meseta 2020--2024};
    \draw[very thick, blue!60!black]
      (0,1.82) -- (2,2.68) -- (4,3.29) -- (6,4.36) -- (8,5.50) -- (8.8,5.55)
      -- (9.2,5.45) -- (9.6,5.48);
    \foreach \x/\y in {0/1.82, 2/2.68, 4/3.29, 6/4.36, 8/5.50, 9.6/5.48} {
      \fill[blue!60!black] (\x,\y) circle (1.6pt);
    }
    \node[font=\scriptsize, above right=0.4mm and 0.8mm, blue!60!black]
      at (0,1.82) {91};
    \node[font=\scriptsize, below right=0mm and 0.8mm, blue!60!black]
      at (4,3.29) {164};
    \node[font=\scriptsize, above=1mm, blue!60!black] at (8,5.50) {275};
    \node[font=\scriptsize, above right=0.8mm and -1mm, blue!60!black]
      at (9.6,5.48) {274};
  \end{tikzpicture}
  \caption{Solicitudes internacionales por la vía PCT, en miles y redondeadas:
  \num{90948} en 2000, unas \num{134000} en 2005, \num{164355} en 2010,
  \num{218000} en 2015, \num{274889} en 2020 y \num{273900} en 2024. Triplicarse
  en veinte años es la señal de globalización; el tramo sombreado es lo que la
  propia OMPI llama una meseta de cinco años, con una caída media del
  \SI{0.1}{\percent} anual entre 2020 y 2024 frente al \SI{5.3}{\percent} anual
  de la década previa. Elaboración propia a partir de
  \textcite{wipo2025pct,wipo2001,wipo2006}.}
  \label{fig:pct}
\end{figure}

\section*{Colaboración y control no son la misma noticia}

Cruzando el lugar de la invención con la titularidad del derecho aparecen cuatro
situaciones distintas (Figura~\ref{fig:cuadrantes}), y el debate mejora cuando
se dice en cuál se está.

\begin{figure}[ht]
  \centering
  \begin{tikzpicture}[
    celda/.style={draw, rounded corners, align=center,
      text width=0.325\textwidth, minimum height=13mm, font=\small}
  ]
    \node[celda, fill=blue!7]   (a) at (0,0)    {\textbf{Explotación global}\\
      \scriptsize Invención nacional protegida en mercados externos};
    \node[celda, fill=green!12] (b) at (6.4,0)  {\textbf{Colaboración global}\\
      \scriptsize Co-invención y co-titularidad entre países};
    \node[celda, fill=gray!12]  (c) at (0,-2.1) {\textbf{Cesión del derecho}\\
      \scriptsize Invención local, titular extranjero por compra o contrato};
    \node[celda, fill=orange!14](d) at (6.4,-2.1){\textbf{Generación global}\\
      \scriptsize I+D deslocalizada, titularidad en la matriz};
    \draw[-{Stealth[length=2mm]}, thick] (-3.0,-3.1) -- (9.2,-3.1)
      node[below left=0.5mm and 0mm, font=\scriptsize]
      {Invención: nacional $\rightarrow$ internacional};
    \draw[-{Stealth[length=2mm]}, thick] (-3.0,-3.1) -- (-3.0,0.9);
    \node[font=\scriptsize, align=center, rotate=90] at (-3.45,-1.1)
      {Titularidad: extranjera $\rightarrow$ local o compartida};
  \end{tikzpicture}
  \caption{Cuatro lecturas posibles de una misma estadística de patentes, según
  la taxonomía de \textcite{archibugi1995}. Solo el cuadrante superior derecho
  describe colaboración en sentido estricto. Elaboración propia.}
  \label{fig:cuadrantes}
\end{figure}

Mi lectura es optimista en el eje horizontal y prudente en el vertical. La
co-invención internacional es una buena noticia casi sin matices: obliga a
divulgar, crea vínculos entre laboratorios y acerca a los países pequeños a
fronteras técnicas que no alcanzarían solos; además, un marco de propiedad
intelectual creíble aumenta la transferencia efectiva de tecnología dentro de
las multinacionales, y no solo el papeleo \parencite{branstetter2006}. El eje
vertical es más delicado: un país puede exhibir excelentes cifras de invención
local y ver, al mismo tiempo, cómo la titularidad migra al exterior. Ahí la
estadística celebra una integración que, vista desde dentro, se parece más a una
maquila de conocimiento.

\section*{Mi posición: un mal necesario}

Creo que la patente es un mal necesario. No discuto sus beneficios, que son
reales: obliga a publicar lo que de otro modo sería secreto industrial, da un
plazo defendible para recuperar la inversión y crea un activo transable que hace
posibles las licencias y las alianzas, es decir, buena parte de la colaboración
internacional que este foro celebra. Lo que no consigo justificar es \textbf{el
costo del proceso, en dinero y en tiempo}.

Ese costo no es una impresión personal. Proteger una invención en los mercados
que importan obliga a repetir el trámite jurisdicción por jurisdicción, con
traducciones, agentes locales y anualidades: la estimación clásica sitúa una
patente europea validada en una decena de países varias veces por encima del
costo de su equivalente estadounidense \parencite{vanpottelsberghe2009}. Con el
tiempo pasa lo mismo: el PCT compra treinta meses de plazo, pero la concesión
suele llegar entre dos y cinco años después del depósito, un horizonte posterior
al ciclo de vida comercial de casi cualquier producto de software. Y lo que se
compra es incierto: los límites del derecho a menudo solo se conocen tras
litigar \parencite{bessen2008,jaffe2004}, la superposición de derechos genera
marañas que encarecen innovar \parencite{shapiro2001,heller1998}, las empresas
manufactureras declaran que fuera de farmacia y química protegen mejor con el
secreto y la ventaja de ser primeras \parencite{cohen2000}, y la evidencia
histórica sugiere que los sistemas de patentes influyen más en la
\emph{dirección} de la innovación que en su cantidad \parencite{moser2013}.

Digo \enquote{mal necesario} y no \enquote{mal} a secas porque es el instrumento
que tenemos. No conozco un mecanismo con alcance global que cumpla a la vez las
tres funciones de divulgar, recompensar y permitir transar. Los sustitutos
resuelven una parte y empeoran otra: el secreto industrial elimina el costo del
trámite y también la divulgación, que es justamente el bien público que la
patente produce; los premios y la compra pública funcionan cuando el Estado sabe
por adelantado qué quiere comprar, y casi nunca lo sabe.

De ahí mi aporte principal al debate, que es la parte que veo menos discutida:
\textbf{la globalización multiplica el defecto que me molesta}. Un sistema caro
se vuelve mucho más caro cuando la protección relevante deja de ser nacional, y
ese costo actúa como un filtro que no selecciona por calidad inventiva sino por
capacidad de pago. Por eso las reformas que reducen fricción, como la patente
unitaria europea vigente desde junio de 2023 \parencite{epo2023}, me parecen más
valiosas que cualquier endurecimiento adicional de los derechos.

\begin{table}[ht]
  \centering
  \small
  \caption{Balance de la internacionalización de las patentes, por dimensión.}
  \label{tab:balance}
  \begin{tabularx}{\textwidth}{@{}>{\raggedright\arraybackslash}p{0.15\textwidth}
    >{\raggedright\arraybackslash}X >{\raggedright\arraybackslash}X@{}}
    \toprule
    \textbf{Dimensión} & \textbf{Ventaja} & \textbf{Inconveniente} \\
    \midrule
    Divulgación
      & Publica la enseñanza técnica en todos los mercados donde se protege
      & Redacción defensiva y reivindicaciones amplias dificultan replicar lo
        divulgado \\
    Colaboración
      & La co-titularidad hace contratables las alianzas entre laboratorios
      & Negociar quién es titular de qué retrasa el proyecto más que la propia
        investigación \\
    Acceso a mercados
      & Una pyme puede licenciar en economías donde nunca operará
      & El derecho vale lo que valga su ejecución judicial local, muy desigual \\
    Costo y tiempo
      & El PCT unifica el depósito y aplaza la decisión treinta meses
      & Multiplica traducciones, agentes y anualidades; la concesión llega tarde
        para ciclos cortos \\
    Control
      & Atrae I+D de multinacionales hacia centros locales
      & La titularidad tiende a concentrarse fuera del país donde se inventó \\
    Equidad
      & Las licencias obligatorias y los consorcios ofrecen válvulas de escape
      & Se activan tarde y con alto costo político, como en el debate sobre los
        ADPIC de 2022 \parencite{wto2022} \\
    \bottomrule
  \end{tabularx}
\end{table}

Cierro con dos preguntas para el grupo. Primera: antes de afirmar que un país se
\enquote{globalizó tecnológicamente}, ¿en qué cuadrante de la
Figura~\ref{fig:cuadrantes} está, si el mismo aumento de solicitudes puede
significar cooperación o transferencia de rentas al exterior? Segunda: si la
tendencia ya no es monótona, y la OMPI misma se pregunta en el título de su
revisión anual si lo que muestra la Figura~\ref{fig:pct} es una meseta o solo
una pausa \parencite{wipo2025pct}, ¿la discusión útil sigue siendo \enquote{patentar o no patentar}, o pasa a ser
cuándo el secreto, la velocidad de ejecución o la publicación defensiva rinden
más que una cartera internacional que se paga en efectivo durante veinte años?

\printbibliography[title={Referencias bibliográficas}]

\end{document}
